\section{Анализ предметной области}

\subsection{Контроль состояния дорожного покрытия}

Состояние дорожного покрытия является одним из ключевых факторов, влияющих на безопасность дорожного движения, комфорт эксплуатации транспортных средств и эффективность транспортной инфраструктуры в целом. Нарушения целостности покрытия приводят к увеличению аварийности, ускоренному износу автомобилей, повышению затрат на ремонт и ухудшению транспортной доступности.

В процессе укладки и эксплуатации дорожное покрытие подвергается постоянному воздействию различных факторов, среди которых:

\begin{itemize}
	\item интенсивная транспортная нагрузка;
	\item температурные перепады;
	\item воздействие атмосферных осадков;
	\item циклы замерзания и оттаивания;
	\item естественное старение материалов;
	\item нарушение технологии строительства и ремонта.
\end{itemize}

Под дефектами дорожного покрытия понимаются локальные или распределённые нарушения структуры полотна, возникающие вследствие механических нагрузок, климатических воздействий, ошибок строительства, износа материалов и недостаточного обслуживания.

К наиболее распространённым дефектам относятся:

\begin{itemize}
	\item выбоины;
	\item продольные и поперечные трещины;
	\item сеточные разрушения;
	\item колейность;
	\item просадки;
	\item шелушение покрытия;
	\item разрушение кромок дорожного полотна.
\end{itemize}

Своевременное обнаружение дефектов позволяет предотвращать дальнейшее разрушение покрытия, снижать вероятность дорожно-транспортных происшествий, оптимизировать планирование ремонтных работ, а также сокращать затраты на содержание дорожной инфрструктуры.

Традиционный контроль состояния дорожного полотна осуществляется посредством визуального обследования специалистами. Такой подход обладает рядом недостатков:

\begin{itemize}
	\item высокая трудоёмкость;
	\item значительные временные затраты;
	\item зависимость от квалификации инспектора;
	\item субъективность оценки;
	\item сложность регулярного мониторинга больших территорий.
\end{itemize}

В связи с этим актуальной задачей является автоматизация процессов диагностики дорожного покрытия с использованием современных методов обработки изображений и машинного обучения.

\subsection{Методы обнаружения дефектов дорожного покрытия}

Существующие методы анализа состояния дорожного полотна можно разделить на несколько основных групп.

\subsubsection{Ручной визуальный контроль}

Данный метод предполагает проведение обследования дорожного покрытия специалистами непосредственно на участке дороги.

Преимуществами метода являются:

\begin{itemize}
	\item высокая точность при участии квалифицированного эксперта;
	\item возможность комплексной оценки состояния покрытия;
	\item отсутствие необходимости в сложном оборудовании.
\end{itemize}

Недостатки метода:

\begin{itemize}
	\item низкая скорость обследования;
	\item высокая стоимость при больших объёмах;
	\item человеческий фактор;
	\item невозможность непрерывного мониторинга.
\end{itemize}

\subsubsection{Контроль с использованием специализированных диагностических комплексов}

Данный подход основан на применении мобильных лабораторий, оснащённых:

\begin{itemize}
	\item лазерными сканерами;
	\item стереокамерами;
	\item профилографами;
	\item датчиками вибрации;
	\item GPS/ГЛОНАСС системами.
\end{itemize}

Такие комплексы обеспечивают высокую точность измерений, однако обладают существенными недостатками:

\begin{itemize}
	\item высокая стоимость оборудования;
	\item сложность эксплуатации;
	\item ограниченная доступность для организаций.
\end{itemize}

\subsubsection{Анализ цифровых изображений}

Метод основан на обработке фотографий дорожного покрытия с целью выявления характерных признаков дефектов.

Источниками изображений могут выступать:

\begin{itemize}
	\item стационарные камеры;
	\item мобильные устройства;
	\item автомобильные видеорегистраторы;
	\item беспилотные летательные аппараты;
	\item дорожные инспекционные комплексы.
\end{itemize}

Преимущества метода:

\begin{itemize}
	\item низкая стоимость внедрения;
	\item возможность масштабирования;
	\item автоматизация анализа;
	\item возможность интеграции в существующие системы мониторинга.
\end{itemize}

\subsection{Методы анализа изображений}

\subsubsection{Компьютерное зрение}

Компьютерное зрение представляет собой направление искусственного интеллекта, связанное с разработкой методов автоматического анализа и интерпретации визуальной информации. Основная цель данной области заключается в создании систем, способных извлекать значимые сведения из изображений и видеопотока аналогично тому, как это делает человек.

Формирование компьютерного зрения как самостоятельной научной дисциплины началось в 1960--1970-х годах. Исследования были ориентированы на решение сравнительно простых задач: выделение контуров объектов, анализ геометрических форм и распознавание символов. Для анализа использовались алгоритмы фильтрации, бинаризации, выделения границ и морфологических преобразований. Среди наиболее известных подходов данного периода можно выделить оператор Собеля, оператор Кэнни и преобразование Хафа.

Постепенно компьютерное зрение развивалось благодаря росту вычислительных мощностей и появлению методов анализа текстур, статистического распознавания образов и первых алгоритмов машинного обучения.

Основными задачами компьютерного зрения являются:

\begin{itemize}
	\item классификация изображений;
	\item сегментация;
	\item обнружение объектов;
	\item отслеживание объектов;
	\item анализ пространственных характеристик сцены.
\end{itemize}

Тем не менее классические методы имеют существенные ограничения. Они требуют ручного проектирования признаков и плохо адаптируются к изменению условий съёмки, освещения, масштаба объектов и наличию шумов.

\subsubsection{Свёрточные нейронные сети}

Значительный этап развития компьютерного зрения связан с появлением искусственных нейронных сетей, а особенно --- свёрточных нейронных сетей (Convolutional Neural Networks, CNN). Их развитие стало возможным благодаря увеличению вычислительных мощностей графических процессоров и накоплению больших размеченных наборов изображений.

Первые идеи использования локальных свёрток для анализа изображений были предложены ещё в 1980-х годах. Существенный прорыв произошёл в 1998 году с появлением архитектуры LeNet, разработанной Яном Лекуном для распознавания рукописных цифр. Действительно широкое распространение свёрточные сети получили после публикации модели AlexNet в 2012 году, которая продемонстрировала значительное превосходство над классическими методами на конкурсе ImageNet.

Основное отличие свёрточных сетей от традиционных алгоритмов заключается в способности автоматически извлекать признаки из изображений: CNN самостоятельно обучается находить информативные признаки на разных уровнях абстракции.

Архитектура свёрточной нейронной сети обычно включает последовательность свёрточных слоёв, функций активации, слоёв подвыборки и полносвязных слоёв. На начальных этапах обработки сеть выявляет простые элементы изображения --- линии, углы, контуры. Более глубокие слои формируют представления сложных структур и объектов.

Для задач анализа дорожного покрытия использование свёрточных нейронных сетей является особенно эффективным, поскольку дефекты асфальтового полотна обладают сложной визуальной структурой и могут значительно отличаться по форме, размеру и условиям съёмки.

Основная цель нейросетевых алгоритмов обнаружения состоит в определении:

\begin{itemize}
	\item принадлежности объекта к определённому классу;
	\item координат ограничивающей рамки объекта;
	\item вероятности корректности предсказания.
\end{itemize}

Современные алгоритмы делятся на две основные категории.

\subsubsection{Двухэтапные детекторы}

К ним относятся семейства:

\begin{itemize}
	\item R-CNN;
	\item Fast R-CNN;
	\item Faster R-CNN;
	\item Mask R-CNN.
\end{itemize}

Принцип работы включает генерацию кандидатов областей и классификацию и уточнение координат.

Преимущества:

\begin{itemize}
	\item высокая точность;
	\item хорошее качество локализации.
\end{itemize}

Недостатки:

\begin{itemize}
	\item высокая вычислительная сложность;
	\item низкая скорость обработки;
	\item ограниченная пригодность для работы в реальном времени.
\end{itemize}

\subsubsection{Одноэтапные детекторы}

К данной категории относятся:

\begin{itemize}
	\item SSD;
	\item RetinaNet;
	\item семейство YOLO.
\end{itemize}

Они выполняют обнаружение объектов за один проход нейронной сети.

Преимущества:

\begin{itemize}
	\item высокая скорость;
	\item возможность работы в режиме реального времени;
	\item хорошее соотношение точности и производительности.
\end{itemize}

Для задачи обнаружения дефектов дорожного покрытия одноэтапные детекторы являются наиболее предпочтительными.

\subsubsection{Сегментационные методы анализа изображений}

Помимо алгоритмов обнаружения объектов, в компьютерном зрении активно применяются методы семантической и instance-сегментации.

Семантическая сегментация предполагает классификацию каждого пикселя изображения, что позволяет получить точное выделение областей, относящихся к определённым классам. Для решения подобных задач широко используются архитектуры:

\begin{itemize}
	\item U-Net;
	\item SegNet;
	\item DeepLab;
	\item PSPNet.
\end{itemize}

Преимуществами сегментационных моделей являются:

\begin{itemize}
	\item высокая точность пространственной локализации;
	\item возможность анализа сложных геометрических структур;
	\item получение детализированных масок объектов;
	\item высокая информативность результатов.
\end{itemize}

Однако применение таких методов связано с рядом ограничений.

Во-первых, для обучения требуется детальная разметка данных, при которой необходимо выделять контуры объектов на уровне пикселей. Подготовка подобных наборов данных требует значительных временных затрат.

Во-вторых, сегментационные модели предъявляют повышенные требования к вычислительным ресурсам и, как правило, уступают детекторам объектов по скорости обработки.

В задачах, где приоритетом является оперативность анализа и возможность работы в режиме реального времени, более рациональным является использование детекторов объектов.

\subsection{Архитектура YOLO}

Семейство алгоритмов YOLO (You Only Look Once) является одним из наиболее известных решений в области обнаружения объектов на текущий момент. Первая версия модели была представлена Джозефом Редмоном в 2015 году и стала значимым этапом развития компьютерного зрения.

Главная идея YOLO заключается в отказе от многоэтапной обработки изображения. Вместо последовательного поиска областей интереса и последующей классификации сеть рассматривает изображение как единую задачу регрессии. Изображение разбивается на сетку, для каждой ячейки которой вычисляются параметры ограничивающих рамок и вероятности присутствия объектов. Такой подход существенно ускоряет обработку по сравнению с двухэтапными детекторами.

Основными преимуществами архитектуры являются:

\begin{itemize}
	\item высокая скорость работы;
	\item компактность моделей, и, как следствие, возможность применения на пользовательских устройствах;
	\item удобство обучения на пользовательских наборах данных;
	\item высокая точность обнаружения при должном обучении.
\end{itemize}

\subsection{Особенности обучения нейросетевых моделей}

Эффективность современных систем компьютерного зрения напрямую зависит от качества обучающих данных.

Процесс подготовки набора данных представляет собой один из наиболее трудоёмких этапов разработки и включает:

\begin{itemize}
	\item сбор исходных изображений;
	\item разметку объектов;
	\item проверку корректности аннотаций;
	\item разделение выборки на обучающую, валидационную и тестовую;
	\item предварительную обработку данных.
\end{itemize}

Одним из важнейших этапов является аугментация данных, позволяющая искусственно расширить обучающую выборку и повысить устойчивость модели к различным условиям.

К основным методам аугментации относятся:

\begin{itemize}
	\item изменение яркости;
	\item изменение контрастности;
	\item поворот изображения;
	\item масштабирование;
	\item зеркальное отражение;
	\item случайное кадрирование;
	\item добавление шумов;
	\item размытие;
	\item изменение цветового баланса.
\end{itemize}

Использование аугментации позволяет моделировать реальные условия эксплуатации системы и снижает вероятность переобучения.

Существенной проблемой также является дисбаланс классов, при котором отдельные категории объектов представлены в выборке значительно чаще других. В подобных условиях требуется дополнительная балансировка обучающих данных.

Качество обучающего набора оказывает непосредственное влияние на точность и устойчивость итоговой модели.

\subsection{Обзор существующих решений}

В России идёт развитие различных интеллектуальных систем мониторинга дорожного покрытия.

Одним из наиболее известных отечественных решений является система <<СКАНДОР>>, предназначенная для автоматизированного выявления дефектов дорожного покрытия и элементов дорожной инфраструктуры с применением технологий искусственного интеллекта. Система ориентирована на использование в составе интеллектуальных транспортных систем и предусматривает интеграцию с муниципальными службами.

Несмотря на высокий технологический уровень подобных решений, большинство существующих систем ориентированы на промышленное или государственное применение и обладают рядом ограничений: высокой стоимостью внедрения, необходимостью специализированного оборудования и закрытостью архитектуры.

\subsection{Выводы по анализу предметной области}

Проведённый анализ показал, что традиционные методы контроля состояния дорожного покрытия не обеспечивают необходимой оперативности и масштабируемости.

Наиболее перспективным направлением является использование методов, основанных на свёрточных нейронных сетях.

Среди современных архитектур обнаружения объектов оптимальным решением для поставленной задачи являются модели семейства YOLO, обеспечивающие баланс между точностью распознавания и скоростью обработки.
